\documentclass[11pt]{article}
\usepackage{mathblatt}
% gesetzt von werkzeuge/setzer.py (Sorte selbst) aus dem Lernweg 4CR
% und den Bankzeilen; Muster bau/proben/2026-10-09/hypotenuse-muster.
\usepackage{enumitem}
\usepackage{adjustbox,varwidth}
\newgeometry{a4paper,margin=18mm,top=15mm,bottom=20mm,footskip=9mm}
\pagestyle{fancy}\fancyhf{}\renewcommand{\headrulewidth}{0pt}
\fancyfoot[L]{\footnotesize\color{mbgrau}4CR \textperiodcentered{} Lösungen \textperiodcentered{} Kenngrößen}
\fancyfoot[R]{\footnotesize\color{mbgrau}\thepage}
\setlength{\parskip}{3pt}
\renewcommand{\kreuz}[1]{\mbox{$\square$\ #1}\quad}
\newcommand{\kreuzl}[1]{\par\hangindent1.4em\hangafter1\noindent$\square$\ #1}
\newcommand{\aufg}[3]{\par\noindent\makebox[0pt][r]{#3}\begin{minipage}[t]{\linewidth}#1\end{minipage}\par}
\newcommand{\aufgg}[4]{\par\noindent\makebox[0pt][r]{#4}\begin{minipage}[t]{0.6\linewidth}\vspace{0pt}#1\end{minipage}\hfill
  \begin{minipage}[t]{0.37\linewidth}\vspace{0pt}\centering #2\end{minipage}\par}
\newcommand{\nr}[1]{\makebox[7mm][l]{\textbf{#1}}}
\newcommand{\lsg}[3]{\par\noindent\begin{minipage}[t]{9mm}\textbf{#1}\end{minipage}%
  \begin{minipage}[t]{52mm}\raggedright\bfseries\boldmath #2\end{minipage}\hspace{3mm}%
  \begin{minipage}[t]{\dimexpr\linewidth-64mm\relax}\raggedright\small #3\end{minipage}\par\vspace{4pt}}
\newlength{\kr}
\makeatletter
\newcommand{\karomerk}[2]{\protected@write\@auxout{}{\string\karoseite{#1}{#2}{\thepage}}}
\newcommand{\karoseite}[3]{\expandafter\xdef\csname karo@#1@#2\endcsname{#3}}
\newcommand{\karohinweis}[3]{\karomerk{h}{#1}\ifcsname karo@k@#1\endcsname
  \ifnum\csname karo@k@#1\endcsname=0\csname karo@h@#1\endcsname\relax#2\else#3\fi
  \else#3\fi}
\makeatother
\begin{document}

\noindent{\Large\bfseries Lösungen}\hspace{1em}{\color{mbgrau}Kenngrößen}\par\smallskip
\noindent{\small Links das Ergebnis, rechts der Weg in kurzen Schritten. Stimmt dein Ergebnis nicht, such den ersten Schritt, der anders ist.}\par
\par\Needspace{25mm}\vspace{6pt}\noindent\colorbox{black!12}{\makebox[7mm]{\bfseries\strut A}}\hspace{2mm}{\bfseries Spannweite und Modalwert}\par\vspace{4pt}
\lsg{1.}{Min.\ $15$, Max.\ $45$, Spannweite $30$\,min}{Minimum $15$ min, Maximum $45$ min; Spannweite $= 45 - 15 = 30$ min}
\lsg{2.}{Spannweite $0{,}20$\,€, Modalwert $0{,}50$\,€}{Spannweite $= 0{,}60 - 0{,}40 = 0{,}20$ €; Modalwert $= 0{,}50$ € (zweimal)}
\lsg{3.}{Super ($0{,}09$\,€ gegen $0{,}06$\,€)}{Super: Spannweite $= 1{,}85 - 1{,}76 = 0{,}09$ €; Diesel: $1{,}72 - 1{,}66 = 0{,}06$ €. Der Preis für Super schwankte stärker.}
\lsg{4.}{Aussage 2; richtig: Modalwert $9$ Jahre}{Modalwert $= 9$ Jahre (dreimal), $13$ ist das Maximum; Aussage 1 stimmt: $13 - 9 = 4$}
\par\Needspace{25mm}\vspace{6pt}\noindent\colorbox{black!12}{\makebox[7mm]{\bfseries\strut B}}\hspace{2mm}{\bfseries Median}\par\vspace{4pt}
\lsg{1.}{$6$ Treffer}{geordnet $3, 4, 6, 7, 9$: Median $= 6$ Treffer}
\lsg{2.}{$13$\,cm}{geordnet $9, 11, 12, 14, 16, 20$: Median $= (12 + 14) : 2 = 13$ cm}
\lsg{3.}{$19{,}5$\,°C}{geordnet $17, 18, 19, 20, 23, 25$: Median $= (19 + 20) : 2 = 19{,}5$ °C}
\lsg{4.}{Ali ($14{,}5$ gegen $14$)}{Tim: geordnet $9, 10, 12, 14, 15, 18, 21$, Median $= 14$; Ali: geordnet $8, 11, 13, 16, 17, 22$, Median $= (13 + 16) : 2 = 14{,}5$. Ali liegt vorn.}
\par\Needspace{25mm}\vspace{6pt}\noindent\colorbox{black!12}{\makebox[7mm]{\bfseries\strut C}}\hspace{2mm}{\bfseries Arithmetisches Mittel}\par\vspace{4pt}
\lsg{1.}{$2$ Tore pro Spiel}{Summe $= 10$; $10 : 5 = 2$ Tore pro Spiel}
\lsg{2.}{$4{,}11$\,m}{Summe $= 16{,}44$ m; $16{,}44 : 4 = 4{,}11$ m}
\lsg{3.}{$\approx 24\,258$ Zuschauer}{$412\,390 : 17 = 24\,258{,}2\ldots \approx 24\,258$ Zuschauer pro Spiel}
\lsg{4.}{Nein, nur $\approx 193$ Besucher am Tag}{Nein: ein Kästchen $= 50$ Besucher; Summe $= 150 + 100 + 150 + 150 + 200 + 300 + 300 = 1\,350$; $1\,350 : 7 = 192{,}8\ldots \approx 193$ Besucher (Personen ganz), also weniger als $200$.}
\par\Needspace{25mm}\vspace{6pt}\noindent\colorbox{black!12}{\makebox[7mm]{\bfseries\strut D}}\hspace{2mm}{\bfseries Mittel rückwärts und Änderungen}\par\vspace{4pt}
\lsg{1.}{$10$}{Summe $= 8 \cdot 3 = 24$; dritte Zahl $= 24 - 5 - 9 = 10$}
\lsg{2.}{$64$\,s}{Summe $= 62 \cdot 4 = 248$ s; bekannt $60 + 65 + 59 = 184$ s; vierte Runde $= 248 - 184 = 64$ s}
\lsg{3.}{$11$ Jahre, es sinkt}{Summe $= 12 \cdot 5 = 60$; neue Summe $= 60 - 16 = 44$; $44 : 4 = 11$ Jahre; das Mittel sinkt}
\lsg{4.}{Ja, mit mindestens $17$ Punkten}{Ja: Summe für fünf Tests $= 15 \cdot 5 = 75$; bisher $14 + 17 + 12 + 15 = 58$; nötig $75 - 58 = 17$ Punkte, das ist höchstens $20$.}
\par\Needspace{25mm}\vspace{6pt}\noindent\colorbox{black!12}{\makebox[7mm]{\bfseries\strut T}}\hspace{2mm}{\bfseries Probetest}\par\vspace{4pt}
\lsg{1.}{$20$\,°C}{geordnet $16, 17, 19, 21, 22, 24$: Median $= (19 + 21) : 2 = 20$ °C}
\lsg{2.}{Do: $95$ Besucher}{Summe $= 120 \cdot 6 = 720$; bekannt $95 + 130 + 110 + 140 + 150 = 625$; Donnerstag $= 720 - 625 = 95$ Besucher}
\lsg{3.}{Aussage 2; richtig: Median $40$}{geordnet $37, 38, 38, 39, 40, 40, 40, 41, 42$, der Median ist der $5.$ Wert: Median $= 40$; Aussage 1 stimmt: $42 - 37 = 5$}
\lsg{4.}{\mbox{a) $46$ Jahre} \mbox{b) Nein, bleibt $46$ Jahre}}{a) Summe $= 506$; $506 : 11 = 46$ Jahre; b) Nein: $506 - 68 - 24 = 414$; $414 : 9 = 46$ Jahre. Die beiden sind zusammen $92 = 2 \cdot 46$ Jahre alt.}
\end{document}
